\chapter{QR and Cholesky Factor Derivatives}
\label{ch:bf-factor-derivatives}

Square-root filters propagate factors rather than covariance matrices.  Their
derivative contract is therefore not only a Kalman contract; it is also a factor
calculus contract.  The likelihood score and Hessian in
Chapters~\ref{ch:bf-kalman-score}--\ref{ch:bf-kalman-hessian} are expressed in
terms of $v_t$, $S_t$, and their derivatives.  A factor backend is acceptable
only if the factors it propagates reproduce those covariance objects and the
same derivative objects.

BayesFilter uses three levels of factor evidence:
\begin{enumerate}[label=(\roman*)]
  \item Cholesky or QR identities stated in the source derivation;
  \item reconstruction diagnostics proving that differentiated factors
    reconstruct the differentiated covariance objects;
  \item backend parity against solve-form likelihood, score, and Hessian.
\end{enumerate}

The QR square-root tests in MacroFinance require prediction, innovation, and
filtered factor derivative reconstruction errors below tight tolerances on
small-dimensional parameter grids, and they compare QR likelihood, gradient,
and Hessian against the solve-form backend.  Those tests are currently the
right standard for BayesFilter factor-derivative chapters: a factor formula is
not production-ready until both reconstruction and end-to-end likelihood
parity pass.

\section{Reconstruction Contract}
\label{sec:bf-factor-reconstruction-contract}

Let $A(\theta)$ be any covariance-like symmetric positive semidefinite object
used by the filter, and let an implemented factor satisfy
\[
  A_\star(\theta)=C(\theta)C(\theta)^\top .
\]
The subscript $\star$ is deliberate.  For a Cholesky backend,
$A_\star=A$ on a positive-definite branch.  For a spectral backend,
$A_\star$ may be a rank-truncated, floored, or symmetrized version of $A$.
The first derivative reconstructed by the factor is
\begin{equation}
\label{eq:bf-factor-reconstruction-first}
  \dot A_\star^{(i)}
  =
  \dot C^{(i)}C^\top
  +C\dot C^{(i)\top}.
\end{equation}
The second derivative reconstructed by the factor is
\begin{equation}
\label{eq:bf-factor-reconstruction-second}
  \ddot A_\star^{(ij)}
  =
  \ddot C^{(ij)}C^\top
  +C\ddot C^{(ij)\top}
  +\dot C^{(i)}\dot C^{(j)\top}
  +\dot C^{(j)}\dot C^{(i)\top}.
\end{equation}
Equations~\eqref{eq:bf-factor-reconstruction-first}--\eqref{eq:bf-factor-reconstruction-second}
are the universal audit equations for every factor backend.  The derivative may
be obtained by analytic formulas, custom gradients, or autodiff, but the
resulting $C$, $\dot C$, and $\ddot C$ must reconstruct the implemented
covariance branch that produced the likelihood value.

\section{Cholesky Factor Derivatives}
\label{sec:bf-factor-cholesky}

Suppose $A=LL^\top$, where $A$ is positive definite and $L$ is lower triangular
with positive diagonal.  Define the lower-triangular half-diagonal operator
\[
  \Phi(X)_{ab} =
  \begin{cases}
    X_{ab}, & a>b,\\
    X_{aa}/2, & a=b,\\
    0, & a<b.
  \end{cases}
\]
For a first derivative, form
\[
  B^{(i)}=L^{-1}\dot A^{(i)}L^{-\top},\qquad
  G^{(i)}=\Phi\!\left(B^{(i)}\right).
\]
Then the Cholesky factor derivative is
\begin{equation}
\label{eq:bf-factor-cholesky-first}
  \dot L^{(i)} = L G^{(i)}.
\end{equation}

For a second derivative, remove the known first-order products before applying
the same triangular split:
\[
  C^{(ij)}
  =
  L^{-1}\!\left(
    \ddot A^{(ij)}
    -\dot L^{(i)}\dot L^{(j)\top}
    -\dot L^{(j)}\dot L^{(i)\top}
  \right)L^{-\top},
  \qquad
  H^{(ij)}=\Phi\!\left(C^{(ij)}\right).
\]
Then
\begin{equation}
\label{eq:bf-factor-cholesky-second}
  \ddot L^{(ij)} = L H^{(ij)}.
\end{equation}
The formulas are local smooth-branch formulas: they require a positive diagonal and
become numerically fragile near a Cholesky pivot boundary.  A backend that adds
jitter before Cholesky must report whether its derivatives target the original
$A$ or the implemented $A+\delta I$.

\section{QR Factor Derivatives}
\label{sec:bf-factor-qr}

QR square-root filters often factor a stack rather than a covariance matrix.  Let
$M(\theta)\in\mathbb{R}^{m\times n}$ have full column rank, $m\ge n$, and use
the thin convention
\[
  M=QR,\qquad Q^\top Q=I_n,\qquad R\ \hbox{upper triangular},\qquad R_{aa}>0.
\]
The positive diagonal fixes the local sign convention.  Define
\[
  A^{(i)} = Q^\top\dot M^{(i)}R^{-1}.
\]
Let $\Omega^{(i)}$ be the skew-symmetric matrix identified from
\[
  \Omega^{(i)}_{ab}=
  \begin{cases}
    A^{(i)}_{ab}, & a>b,\\
    0, & a=b,\\
    -A^{(i)}_{ba}, & a<b,
  \end{cases}
  \qquad
  T^{(i)}=A^{(i)}-\Omega^{(i)} .
\]
Then
\begin{align}
\label{eq:bf-factor-qr-r-first}
  \dot R^{(i)} &= T^{(i)}R,\\
\label{eq:bf-factor-qr-q-first}
  \dot Q^{(i)} &= Q\Omega^{(i)}+(I-QQ^\top)\dot M^{(i)}R^{-1}.
\end{align}

For second derivatives, define the effective second-order stack derivative
\[
  E^{(ij)}
  =
  \ddot M^{(ij)}
  -\dot Q^{(i)}\dot R^{(j)}
  -\dot Q^{(j)}\dot R^{(i)}
\]
and
\[
  B^{(ij)}=Q^\top E^{(ij)}R^{-1},\qquad
  C_Q^{(ij)}
  =
  -\dot Q^{(i)\top}\dot Q^{(j)}
  -\dot Q^{(j)\top}\dot Q^{(i)}.
\]
The matrix $\Gamma^{(ij)}=Q^\top\ddot Q^{(ij)}$ is determined by
\[
  \Gamma^{(ij)}_{ab}=
  \begin{cases}
    B^{(ij)}_{ab}, & a>b,\\
    C_{Q,aa}^{(ij)}/2, & a=b,\\
    C_{Q,ab}^{(ij)}-B^{(ij)}_{ba}, & a<b,
  \end{cases}
  \qquad
  U^{(ij)}=B^{(ij)}-\Gamma^{(ij)}.
\]
Then
\begin{align}
\label{eq:bf-factor-qr-r-second}
  \ddot R^{(ij)} &= U^{(ij)}R,\\
\label{eq:bf-factor-qr-q-second}
  \ddot Q^{(ij)} &= Q\Gamma^{(ij)}+(I-QQ^\top)E^{(ij)}R^{-1}.
\end{align}
The QR formulas are local to a fixed rank and sign branch.  Pivoting,
column-rank loss, or diagonal sign corrections are branch events and must be
reported as derivative-policy events, not hidden inside a smooth Hessian claim.

\section{Square-Root Kalman Stacks}
\label{sec:bf-factor-kalman-stacks}

The generic formulas above enter the filter through stack factorizations.  If
$C_{t-1|t-1}C_{t-1|t-1}^\top=P_{t-1|t-1}$ and
$G_tG_t^\top=Q_t$, the prediction covariance can be declared by the stack
\[
  M_t^P =
  \begin{bmatrix}
    F_tC_{t-1|t-1} & G_t
  \end{bmatrix},
  \qquad
  P_{t|t-1}=M_t^PM_t^{P\top}.
\]
Equivalently, a thin QR factorization of $M_t^{P\top}=Q_t^PR_t^P$ gives the
lower factor $C_{t|t-1}=R_t^{P\top}$.  The first and second derivatives of
$M_t^P$ are obtained by product rule from the derivatives of
$F_t$, $C_{t-1|t-1}$, and $G_t$, and
\eqref{eq:bf-factor-qr-r-first}--\eqref{eq:bf-factor-qr-r-second} then give the
derivatives of $C_{t|t-1}$.

Likewise, if $E_tE_t^\top=R_t$,
\[
  M_t^S =
  \begin{bmatrix}
    H_tC_{t|t-1} & E_t
  \end{bmatrix},
  \qquad
  S_t=M_t^SM_t^{S\top},
\]
and QR of $M_t^{S\top}$ gives the innovation factor used by the solve-form
likelihood.  The filtered covariance can be audited by the Joseph stack
\[
  M_t^U =
  \begin{bmatrix}
    (I-K_tH_t)C_{t|t-1} & K_tE_t
  \end{bmatrix},
  \qquad
  P_{t|t}=M_t^UM_t^{U\top}.
\]
Derivatives of $K_t$ are those in
\eqref{eq:bf-score-gain-first} and
\eqref{eq:bf-hessian-gain-second}.  Therefore a square-root implementation has
two equivalent audit paths: reconstruct
\eqref{eq:bf-factor-reconstruction-first}--\eqref{eq:bf-factor-reconstruction-second}
from differentiated factors, and compare the resulting score/Hessian objects to
the covariance-form recursions in Chapters~\ref{ch:bf-kalman-score}--\ref{ch:bf-kalman-hessian}.

\section{Spectral Factors}
\label{sec:bf-factor-spectral}

SVD sigma-point filters use a spectral factor, not a triangular factor, to
generate point displacements.  If
\[
  P=U\Lambda U^\top,\qquad C=U\Lambda_+^{1/2},
\]
then the factor actually reconstructs
\[
  P_+ = CC^\top = U\Lambda_+U^\top,
\]
where $\Lambda_+$ includes the rank, floor, or smoothing policy used by the
implementation.  Therefore spectral factor validation must distinguish two
questions:
\begin{enumerate}[label=(\roman*)]
  \item Do $C$ and its derivatives reconstruct $P_+$ and its derivatives?
  \item If $P_+\ne P$, is the difference a declared model approximation, a
    finite-arithmetic regularization, or an invalid fallback?
\end{enumerate}
This distinction is essential for HMC.  A hard eigenvalue floor can produce a
finite likelihood while changing the derivative target at the floor boundary.
If the floor is hard, derivatives are branch derivatives of $P_+$, not
derivatives of the pre-regularized covariance $P$.  If the floor is smooth, the
floor derivative is part of the model actually differentiated.  Either way, the
factor reconstruction equations in
\eqref{eq:bf-factor-reconstruction-first}--\eqref{eq:bf-factor-reconstruction-second}
must be checked against $P_+$.

Chapter~\ref{ch:bf-svd-sigma-point} therefore treats spectral derivatives as a
separate contract with gap telemetry, active-floor reporting, and branch
diagnostics.  Its simple-spectrum SVD/eigen formulas define one smooth-branch
policy; they do not remove the need for finite-difference checks across spectral
gap, rank, and floor regimes.

\section{Backend Parity Gates}
\label{sec:bf-factor-backend-parity}

A factor backend is ready to support analytic gradients and Hessians only after
the following checks pass on controlled systems:
\begin{enumerate}[label=(\roman*)]
  \item value reconstruction:
    $\|A_\star-CC^\top\|$ is within tolerance for prediction, innovation, and
    filtered covariance objects;
  \item first- and second-derivative reconstruction:
    \eqref{eq:bf-factor-reconstruction-first}--\eqref{eq:bf-factor-reconstruction-second}
    match the declared $\dot A_\star$ and $\ddot A_\star$;
  \item solve parity:
    the factor solve gives the same $w_t$, score, and Hessian contribution as
    the covariance solve on well-conditioned systems;
  \item branch reporting:
    jitter, rank truncation, pivoting, sign flips, active floors, and fallback
    paths are emitted as diagnostics;
  \item backend parity:
    Cholesky-refactorized, QR square-root, and SVD/spectral backends agree where
    their mathematical branches represent the same covariance law.
\end{enumerate}
These gates prevent a common mistake: a numerically stable factor can make the
likelihood finite while silently changing the derivative target.  The factor
chapter therefore supplies the algebraic contracts that the validation chapter
turns into executable tests.
